\chapter{Calcium (Blood)}

\section*{What Is This Test?}
The blood calcium test measures the total calcium concentration in the serum or plasma. Calcium is the most abundant mineral in the human body. Approximately 99\% of total body calcium is stored in bone as hydroxyapatite crystals [Ca\textsubscript{10}(PO\textsubscript{4})\textsubscript{6}(OH)\textsubscript{2}], and only 1\% circulates in the extracellular fluid. In serum, total calcium exists in three forms: (1) ionized (free) calcium (about 50\%), which is the physiologically active form and the only form regulated by hormones; (2) protein-bound calcium (about 40\%), mostly bound to albumin; and (3) complexed calcium (about 10\%), bound to anions such as phosphate, citrate, sulfate, and bicarbonate. Calcium homeostasis is tightly regulated by three hormones: parathyroid hormone (PTH), which increases serum calcium by stimulating bone resorption, renal tubular calcium reabsorption, and calcitriol synthesis; 1,25-dihydroxyvitamin D (calcitriol), which increases intestinal calcium absorption; and calcitonin, which modestly lowers serum calcium by inhibiting bone resorption. The total calcium test is the most commonly ordered calcium measurement, but it can be misleading when albumin is abnormal. The ionized calcium test measures the physiologically active fraction directly and is preferred in critical illness, massive transfusion, and when the protein status is abnormal.

\section*{Alternative Names}
\begin{itemize}
  \item Calcium, Total
  \item Serum Calcium
  \item Total Calcium (tCa)
  \item Ionized Calcium (iCa)
  \item Ca, Blood
  \item Corrected Calcium
\end{itemize}
\section*{Principle}
Total serum calcium is measured by automated analyzers using one of several methods. The most common are colorimetric methods using o-cresolphthalein complexone (OCPC) or Arsenazo III dye, which form colored complexes with calcium that absorb at a specific wavelength (570--650 nm). Alternatively, ion-selective electrode (ISE) methods measure the calcium activity in serum or plasma directly, reporting ionized calcium. The normal total calcium range is 8.5--10.5 mg/dL (2.12--2.62 mmol/L). However, because about 40\% of total calcium is bound to albumin, total calcium must be interpreted with the serum albumin. When albumin is abnormal, a corrected calcium is calculated using the formula: Corrected Ca (mg/dL) = Measured Ca (mg/dL) + 0.8 $\times$ (4.0 -- Albumin in g/dL). This formula (the Payne correction) is approximate and does not apply in acid-base disorders or when ionized calcium is preferred. Ionized calcium is the physiologically active form and is directly measured with an ion-selective electrode, usually on a blood-gas analyzer. The normal ionized calcium range is 4.5--5.6 mg/dL (1.13--1.40 mmol/L). Ionized calcium is affected by pH: alkalosis increases protein binding and lowers ionized calcium; acidosis has the opposite effect. The calcium-phosphate product (Ca $\times$ PO\textsubscript{4}) is clinically important: a product >70 mg$^2$/dL$^2$ in chronic kidney disease increases the risk of vascular calcification and calciphylaxis. Calcium is regulated by PTH and calcitriol (1,25-dihydroxyvitamin D): PTH increases calcium by mobilizing bone calcium, increasing renal tubular reabsorption, and stimulating calcitriol; calcitriol increases intestinal absorption. The calcium-sensing receptor (CaSR) on the parathyroid chief cells detects serum ionized calcium and modulates PTH secretion in a classic negative feedback loop. This is the basis of familial hypocalciuric hypercalcemia (FHH), where an inactivating CaSR mutation causes mild hypercalcemia with inappropriately normal PTH and low urine calcium excretion.

\section*{History}
Calcium was identified as an element by Humphry Davy in 1808. The relationship between the parathyroid glands and calcium homeostasis was established in the 1920s--1930s. In 1925, James Collip isolated an active parathyroid extract (PTH), and hyperparathyroidism was recognized as a cause of hypercalcemia. The radioimmunoassay for PTH was developed by Berson and Yalow in the 1960s, which made it possible to measure PTH directly. In the 1990s, the calcium-sensing receptor (CaSR) was cloned by Brown and colleagues, and mutations in it were found to cause familial hypocalciuric hypercalcemia and neonatal severe hyperparathyroidism. The correction formula for albumin-bound calcium was introduced in the 1970s (Payne formula). Modern automated analyzers measure total calcium by colorimetric methods and ionized calcium by ion-selective electrodes. Guidelines for the management of hypercalcemia and hypocalcemia were published by major endocrine societies \cite{bilezikian2018}\cite{cooper2008}\cite{stewart2005}.

\section*{Why This Test Is Ordered}
\begin{itemize}
  \item Evaluation of hypercalcemia: calcium is measured in patients with weakness, fatigue, confusion, constipation, polyuria, nephrolithiasis, or bone pain. Hypercalcemia is most commonly caused by primary hyperparathyroidism (in outpatients) or malignancy (in hospitalized patients)
  \item Evaluation of hypocalcemia: calcium is measured in patients with paresthesias, muscle cramps, tetany, seizures, or prolonged QT interval. Hypocalcemia may be caused by hypoparathyroidism, vitamin D deficiency, renal failure, or hypomagnesemia
  \item Monitoring during therapies that affect calcium: cisplatin, denosumab, bisphosphonates, loop diuretics, glucocorticoids, calcimimetics (cinacalcet), and after parathyroidectomy
  \item Routine chemistry panels (BMP/CMP) and assessment of renal function: calcium is an integral component of the comprehensive metabolic panel, and the calcium-phosphate product is monitored in CKD
  \item Preoperative evaluation: calcium is measured before parathyroid and thyroid surgery and in patients with nephrolithiasis (to exclude hyperparathyroidism)
\end{itemize}

\section*{Primary vs Secondary Test}
Total serum calcium is a primary screening test. Ionized calcium is a secondary test used in specific situations: critical illness, massive blood transfusion (citrate), cardiac surgery, hypoalbuminemia, acid-base disorders, and when total calcium is unreliable.

\section*{How This Test Is Performed}
A healthcare professional draws blood from a vein into a serum separator tube (gold-top) or lithium heparin tube (green-top). The sample is centrifuged and analyzed by colorimetric methods (total calcium) or by an ion-selective electrode (ionized calcium, usually on a blood-gas analyzer). Total calcium is stable and can be measured within 1--2 hours. Ionized calcium requires anaerobic collection (to prevent pH changes from carbon dioxide loss) and rapid analysis (within 30 minutes) in a heparinized syringe.

\section*{What Preparation Is Needed}
No fasting is strictly required, but fasting is often recommended for consistency (a high-calcium meal can transiently affect calcium). The patient should be seated for 5--10 minutes before venipuncture (standing increases albumin and total calcium). Document all medications, especially calcium or vitamin D supplements, diuretics, bisphosphonates, and corticosteroids.

\section*{Contraindications and Precautions}
\begin{itemize}
  \item Total calcium must always be interpreted with serum albumin or corrected for albumin. A low total calcium with low albumin may represent a normal ionized calcium (pseudohypocalcemia). Conversely, an elevated total calcium with elevated albumin (dehydration, multiple myeloma with paraprotein) may not reflect true hypercalcemia. When in doubt, measure ionized calcium
  \item Ionized calcium is pH-dependent: respiratory or metabolic alkalosis decreases ionized calcium (increased protein binding), while acidosis increases it. An anaerobic sample and rapid analysis are essential; a sample exposed to air loses CO\textsubscript{2}, raising pH and falsely lowering ionized calcium
  \item Prolonged tourniquet application (stasis), fist clenching, and hemoconcentration falsely elevate calcium by increasing protein concentration
  \item Hemolysis can falsely alter calcium results in some assays (calcium is released from red cells); hemolyzed samples should be rejected
  \item Many drugs affect calcium: thiazide diuretics, lithium, vitamin D, calcium supplements, and tamoxifen can raise calcium; loop diuretics, corticosteroids, bisphosphonates, cinacalcet, and EDTA can lower it
  \item Paraproteins (multiple myeloma, Waldenstr\"{o}m macroglobulinemia) can artifactually elevate total calcium by binding calcium, even when ionized calcium is normal
  \item In chronic kidney disease (CKD), total calcium may be normal despite abnormal ionized calcium; ionized calcium is preferred for management decisions in this population
\end{itemize}
\section*{Risks}
Standard venipuncture risks: mild pain, bruising, or hematoma at the puncture site.

\section*{What Happens After the Test}
Results are available within 1--2 hours. Total calcium is interpreted with albumin and the clinical context. If calcium is abnormal, the provider evaluates the PTH-calcium-vitamin D axis: PTH, phosphate, vitamin D, magnesium, and urine calcium. A single abnormal value is confirmed with a repeat measurement before initiating treatment. Severe hypercalcemia (>14 mg/dL) or symptomatic hypercalcemia requires urgent treatment (IV fluids, calcitonin, bisphosphonates).

\section*{Understanding Results}

\subsection*{Below Normal}
\begin{itemize}
  \item \textbf{Hypocalcemia (total Ca <8.5 mg/dL, ionized Ca <4.5 mg/dL):} May cause paresthesias, muscle cramps, carpal-pedal spasm, Chvostek and Trousseau signs, seizures, and QT prolongation. Causes include hypoparathyroidism (post-thyroidectomy, autoimmune, DiGeorge syndrome), vitamin D deficiency or resistance, hypomagnesemia (inhibits PTH secretion), CKD (phosphate retention, calcitriol deficiency), acute pancreatitis (calcium sequestration), massive transfusion (citrate chelation), and medications (bisphosphonates, loop diuretics, cisplatin)
  \item \textbf{Pseudohypocalcemia:} Low total calcium with normal ionized calcium due to hypoalbuminemia (cirrhosis, nephrotic syndrome, malnutrition). No treatment is needed; the ionized calcium is normal
  \item \textbf{Corrected calcium in hypoalbuminemia:} The corrected calcium formula (Ca + 0.8 $\times$ (4.0 -- albumin)) provides an estimate of total calcium normalized to a normal albumin; it is unreliable when albumin is very low or in acid-base disorders
  \item \textbf{Low ionized calcium in the setting of alkalosis:} Respiratory alkalosis (hyperventilation) or metabolic alkalosis increases calcium-protein binding and lowers ionized calcium, causing symptoms of hypocalcemia (tetany) even when total calcium is normal. Treatment is correction of the alkalosis
\end{itemize}

\subsection*{Above Normal}
\begin{itemize}
  \item \textbf{Primary hyperparathyroidism:} The most common cause of hypercalcemia in outpatients (approximately 1 in 500 women). Excess PTH from a parathyroid adenoma (80--85\%), hyperplasia (15\%), or carcinoma (1\%) causes hypercalcemia with a high or inappropriately normal PTH. Patients may be asymptomatic or have kidney stones, bone loss, or peptic ulcer disease
  \item \textbf{Malignancy-associated hypercalcemia:} The most common cause of hypercalcemia in hospitalized patients. Mechanisms include PTHrP secretion (squamous cell lung cancer, renal cell cancer, breast cancer), osteolytic bone metastases (breast, myeloma), and calcitriol overproduction (lymphoma). PTH is suppressed
  \item \textbf{Mild hypercalcemia with normal or mildly elevated PTH:} Consider familial hypocalciuric hypercalcemia (FHH; CaSR mutation), which is benign and does not require surgery. FHH is distinguished from primary hyperparathyroidism by a low urine calcium-to-creatinine clearance ratio (<0.01) and a family history of hypercalcemia
  \item \textbf{Thiazide diuretics and lithium:} Thiazides reduce renal calcium excretion and can unmask or cause hypercalcemia. Lithium alters the set point of the CaSR and causes mild hypercalcemia in 10--20\% of users
  \item \textbf{Vitamin D excess and granulomatous disease:} Excessive vitamin D or 1,25-dihydroxyvitamin D from granulomas (sarcoidosis, TB) increases intestinal calcium absorption. Milk-alkali syndrome (excess calcium and antacid intake) is an important cause in the elderly
\end{itemize}

\section*{Normal Results}
\begin{itemize}
  \item \textbf{Total calcium (adults):} 8.5--10.5 mg/dL (2.12--2.62 mmol/L)
  \item \textbf{Ionized calcium (adults):} 4.5--5.6 mg/dL (1.13--1.40 mmol/L)
  \item \textbf{Total calcium in children:} 8.5--10.5 mg/dL (slightly higher during growth)
  \item \textbf{Corrected calcium:} Calculated when albumin is abnormal; target normal range 8.5--10.5 mg/dL
\end{itemize}

\section*{Follow-up Tests}
\begin{itemize}
  \item \textbf{Parathyroid hormone (PTH, intact):} The key test for hypercalcemia evaluation. A high or inappropriately normal PTH with hypercalcemia indicates primary hyperparathyroidism. A suppressed PTH (<10 pg/mL) with hypercalcemia indicates non-PTH-mediated causes (malignancy, vitamin D excess, sarcoidosis, milk-alkali syndrome)
  \item \textbf{Phosphate, magnesium, and albumin:} Hypophosphatemia accompanies primary hyperparathyroidism and vitamin D deficiency; hyperphosphatemia accompanies renal failure and hypoparathyroidism. Hypomagnesemia must be corrected or PTH secretion will remain impaired
  \item \textbf{25-hydroxyvitamin D and 1,25-dihydroxyvitamin D:} To evaluate vitamin D status. Low 25-OH vitamin D indicates deficiency; elevated 1,25-(OH)2 vitamin D suggests granulomatous disease or lymphoma
  \item \textbf{24-hour urine calcium:} In hypercalcemia with elevated PTH, a 24-hour urine calcium distinguishes primary hyperparathyroidism (high or normal) from familial hypocalciuric hypercalcemia (low, urine calcium-to-creatinine clearance ratio <0.01). In nephrolithiasis, urine calcium identifies hypercalciuria
  \item \textbf{Ionized calcium:} When total calcium is unreliable (hypoalbuminemia, critical illness, massive transfusion, acid-base disorders). Ionized calcium is the physiologically relevant fraction
  \item \textbf{PTHrP, calcitriol, or imaging:} PTHrP (humoral hypercalcemia of malignancy), 1,25-(OH)2 vitamin D (lymphoma/sarcoidosis), and imaging (sestamibi scan, ultrasound for parathyroid localization; CT for malignancy) as indicated by the clinical picture
\end{itemize}

\section*{References}
\bibliographystyle{unsrtnat}
\bibliography{references}

\clearpage
\section*{Regulatory Status and Authorization Date}
This chapter describes a test category, not a specific commercial product. FDA, CDSCO, EU IVDR, and MHRA approval or registration dates therefore cannot be assigned without the assay manufacturer, product name, instrument, and intended use. For laboratory-developed tests, an individual agency approval date may not exist. See the Preface for the interpretation of regulatory status.

\clearpage
